\section{总结与延伸}

本节把全片压缩成一个可复用框架。2025 年 Q1 的大模型季报可以理解为“智能主线回归”：不要被单个产品爆点或单个公司发布带偏，而要持续追问模型上限从哪里来、模型如何行动、行动如何产生反馈、反馈如何回到学习、产品如何承接智能、公司如何把组织能力转成战略。

\begin{importantbox}{六个核心结论}
第一，Pre-training 仍是底座能力上限的关键，不应因为后训练热度而被忽视。第二，Coding 是模型最早获得可验证行动能力的赛博环境。第三，Agent 是从语言输出到任务执行的系统跃迁。第四，Online Learning 可能把使用过程变成新数据生产过程。第五，应用壁垒来自环境、状态、反馈和成本控制。第六，中美竞争最终要看可验证创造，而不是圈层叙事。
\end{importantbox}

\subsection{本期术语速查表}

本节给出速查，方便把这期和前后张小珺 AI/互联网队列连接起来。EP97 与 EP101 的 Agent 产品、EP102 的多模态路线、EP106/109 的具身智能、EP110 的 Agent 技术报告、EP127/136 的大模型季报有明显连续性。

\noindent\begin{tabularx}{\textwidth}{p{0.23\textwidth}p{0.35\textwidth}X}
\toprule
术语/公司 & 本期含义 & 后续观察方式 \\
\midrule
Pre-training & 打开 base model 内在上限 & 看下一代模型是否出现新能力，而非只刷分。 \\
Reasoning/RL & 强化推理和任务表现 & 看是否能转化为真实任务完成。 \\
Coding & 数字行动环境与模型的手 & 看开发者付费、工具调用和端到端任务成功率。 \\
Agent & 感知、工具、记忆和执行闭环 & 看长任务、错误恢复和可控性。 \\
Online Learning & 使用中产生反馈并持续学习 & 看环境、奖励、记忆和安全审计是否成熟。 \\
DeepSeek & 开源效率和中国模型信号 & 看能力、成本、生态和持续迭代。 \\
Manus/Cursor & 模型能力的产品容器 & 看工作流壁垒、收入和 token 经济。 \\
\bottomrule
\end{tabularx}

\subsection{后续观察问题}

本节把季报转成后续跟踪清单。读者可以用这些问题判断 2025 年后续几个季度，广密这套判断是否继续成立。

\begin{enumerate}
\item 下一代 base model 是否会出现明显新能力，还是主要靠 post-training 刷 benchmark？
\item Coding 场景的领先模型是否继续由 Anthropic/Claude 系列主导，还是 OpenAI、Google、DeepSeek 等会重新追上？
\item Agent 产品是否能从演示走向稳定交付，尤其是在长任务、权限、安全和成本上是否过关？
\item Online Learning 是否会出现可复现的工程范式，而不是停留在概念层？
\item 模型公司与应用公司之间，谁更能掌握用户状态、任务环境和反馈数据？
\item DeepSeek 式开源效率路线是否能持续扩散到更多模型和工具生态？
\item 中美 AI 竞争中，算力限制是否会被算法效率、工程优化和开源生态部分抵消？
\end{enumerate}

\subsection{拓展阅读}

\begin{itemize}
\item 对 Agent 产品和应用创业感兴趣，可对照 EP101 YouWare、EP110 Agent 技术报告、EP139 Agent 技术史。
\item 对多模态和世界模型感兴趣，可对照 EP102 张祥雨访谈、EP133 谢赛宁访谈。
\item 对具身智能和机器人主线感兴趣，可对照 EP106 王鹤、EP109 光轮智能、VLA 投屏版 `eiQFomOuCJs`。
\item 对大模型季报连续性，可对照 EP127、EP136 以及视频描述中列出的 2023/2024 年季度回顾。
\end{itemize}

\end{document}
